% !TeX root = ../main.tex
% =================================================================================
\chapter{Concept}\label{chap:concept}
% =================================================================================
% TODO: Rename this chapter to reflect your approach. Split it into several chapters if
% necessary, e.g., one chapter per research question or per part of your concept.

This chapter presents our concept for \ldots (contribution \refC{concept}), which addresses \refRQ{detect} and \refRQ{resolve}.
\Cref{sec:concept:overview} gives an overview, and \cref{sec:concept:first,sec:concept:second} describe its parts in detail.

% ---------------------------------------------------------------------------------
\section{Overview}\label{sec:concept:overview}
% ---------------------------------------------------------------------------------

Give an overview of your concept before you go into detail, e.g., with a figure such as \cref{fig:concept:overview}.
Reference every figure in the text and explain how to read it:
Do not expect the reader to know what they should see in a figure.

\begin{figure}[htbp]
  \centering
  \begin{tikzpicture}[
      node distance=1.2cm and 1.6cm,
      box/.style={draw=TUMBlue, fill=TUMBlue!8, rounded corners, minimum width=2.6cm, minimum height=1cm, align=center, font=\sffamily\small},
      arrow/.style={-{Stealth}, thick, TUMDarkGray},
    ]
    \node[box] (input) {Input};
    \node[box, right=of input] (approach) {Your\\Approach};
    \node[box, right=of approach] (output) {Output};
    \draw[arrow] (input) -- (approach);
    \draw[arrow] (approach) -- (output);
  \end{tikzpicture}
  \caption[Overview of the concept]{Overview of the concept. A caption must be understandable on its own.}
  \label{fig:concept:overview}
\end{figure}

% ---------------------------------------------------------------------------------
\section{First Part of the Concept}\label{sec:concept:first}
% ---------------------------------------------------------------------------------

Describe the first part of your concept.
Explain your design decisions, the alternatives you considered, and why you chose this option.
Refer to the requirements it fulfills, e.g., FR1 in \cref{tab:requirements:functional}.
Use a running example to illustrate the concept.

% ---------------------------------------------------------------------------------
\section{Second Part of the Concept}\label{sec:concept:second}
% ---------------------------------------------------------------------------------

Proceed in the same way for the further parts of your concept.
